\chapter{An on-chain system for asynchronous transactions}
\label{chap:async-transactions}

\section{Definitions}
We begin the description of our system by defining a few important terms.

An \textbf{user}, denoted using $U_i$, is an entity that interacts with the system to register their asynchronous \textbf{intents of execution} $\Psi_j$. A user can be a person, a smart contract, or a machine. The subscripts $i$ and $j$ are used to distinguish between different users and intents, respectively and one user $U_i$ can have more than one intent registered at the same time.

Consequently, each $\Psi_j$ has a corresponding intent identifier $j$ and has a unique user $i$ who has registered it for asynchronous execution.

Each intent of execution $\Psi_j$ is composed of a list of \textbf{actions} $\{\Psi_{j, 0} ... \Psi_{j, n}\}$ that are to be executed in the specified order within a given \textbf{execution opportunity window (EOW)}. The EOW is a time period during which the actions of the intent can be executed by the system and is specific to an intent, progressing the intent to its completion and yielding fee for the solvers involved. However, the user may want to describe a condition (such as a certain period of time after the first action was executed, the "time to live") under which the intent itself should be \textit{collapsed}, in effect making it impossible to progress further.

\begin{lstlisting}[language=Rust, caption=Intent of Execution]
struct Intent {
    intent_identifier:  u64,
    intent_owner:       Address,
    actions:            Vec<Action>, // count should be available using `actions.len()`
    isEowCollapsed:     Fn(&self) -> bool,
}
\end{lstlisting}

\begin{lstlisting}[language=Rust, caption=Action]
struct Action {
    action_identifier:  u64,
    intra_dependencies: Vec<(u64, u64)>, // list of intra-intent dependencies
    inter_dependencies: Vec<(u64, u64)>, // list of inter-intent dependencies
    accessed_addresses: Vec<Address>, // list of dApps and addresses accessed by the action
}
\end{lstlisting}

\section{Time Turner}
The time turner is a smart contract or a set of smart contracts that allows users to register their intents of execution. The time turner is responsible for managing the registration of intents and ensuring their correct execution. It is also responsible for managing the state of the system and ensuring that all intents are executed in a consistent manner. Also, time turner provides a mechanisms against different actor's misbehaviors and penalizes according to the "asynchronous execution protocol" (AEP) rules. 

Time turner maintains a mapping of all registered intents $\Psi_j$ to their corresponding user identifiers $U_i$, called the \textbf{intent-to-user-map} which exposes the following function: \code{getIntentOwner(j) $\to$ i}.



\section{Asynchronous Execution Protocol (AEP)}

This section describes the Asynchronous Execution Protocol (AEP) that is used to execute the intents of execution registered by users. We first describe the different components of the AEP and then describe the full flow of the protocol at the end.

\subsection{Inter-Intent Dependencies}
An action $\Psi_{j, k}$ can conditionally depend on the successful execution of one or more actions from different intent than $\Psi_{j}$: $\Psi_{x, y}, \Psi_{a, b}, ...$ under intents that are registered by the same user or a different user. This is done by specifying a set of \textbf{inter-intent dependencies} $\{\text{INTER}_{j, k} ... \text{INTER}_{j, n}\}$ for each action in the intent. The inter-intent dependencies are used to ensure that the actions of the intent are executed in the correct order and that the dependencies are satisfied before executing the action.

Each inter-intent dependency $\text{INTER}_{j, k}$ is essentially a list of action references in form of a tuple $(x, y)$ where $x$ is the intent identifier and $y$ is the action identifier within the intent $x$.

For example, if some action $\Psi_{13, 2}$ of intent $\Psi_{13}$ depends on the successful execution of action $\Psi_{12, 5}$ of intent $\Psi_{12}$ and the successful execution of action $\Psi_{7, 1}$ of intent $\Psi_{7}$, then the inter-intent dependency is represented as $\text{INTER}_{13, 2} = [(12, 5), (7, 1)]$. Doing this allows $\Psi_{13, 2}$ to be executed only after some successful execution of both $\Psi_{12, 5}$ and $\Psi_{7, 1}$ and during the execution of $\Psi_{13, 2}$, the time turner will provide the context object generated by the execution of $\Psi_{12, 5}$ and $\Psi_{7, 1}$ to the execution of $\Psi_{13, 2}$.

Two forms of dependencies are possible: 
\begin{itemize}
    \item \textbf{hard dependency}: $\Psi_{j, k}$ executes only if all the dependencies are executed and the intent they are part of has transitioned into the status of \texttt{COMPLETED}.
    \item \textbf{soft dependencies}: $\Psi_{j, k}$ executes only if all the dependencies are executed, without any care of the status of the intent they are part of. These dependencies while executing faster than hard dependencies, are subject to the risk of consuming a wrong action context object (since the action they depend on may have executed differently post execution forking) or even worse, the intent may have failed, leading to all side effects being rolled back and ADRs not being redeemed eventually.
\end{itemize}

\todo{Note: Side channel context injection is possible via solver's runtime inputs. However not enforced by the time turner.}

\subsection{Intra-Intent Dependencies}
As a counternote to "Inter-Intent Dependencies", we also have "Intra-Intent Dependencies". These are dependencies that exist within the same intent. For example, if action $\Psi_{j, k}$ depends on the successful execution of action $\Psi_{j, m}$, then the intra-intent dependency is represented as $\text{INTRA}_{j, k} = (j, m)$. This is used to ensure that the actions of the intent are executed in the correct order and can potentially be parallelized if they do not depend on each other.

It is worth noting that Intra-Intent dependencies are not strictly soft-dependencies only. Hard dependency under intra-intent regime would lead to a deadlock situation, the action cannot execute until the whole intent is executed and the intent cannot be fully executed until the constituent actions are all executed. Hence, the time turner does not allow hard dependencies under intra-intent regime.

\subsection{Anatomy of an Action}
Each actions $\Psi_{j, k}$ of an intent $\Psi_j$ is composed of three components:
\begin{itemize}
    \item \textbf{Action Pre-condition}: A unique identifier for the action within the intent. It is denoted as $k$.
    \item \textbf{Action Execution}: The entity responsible for executing the action. It is denoted as $S_{j, k}$.
    \item \textbf{Action Context update}: The receipt generated by the action solver upon successful execution of the action. It is denoted as $\Delta_{j, k}$.
\end{itemize}

\subsection{Execution Forking for optimal fee}

Assume an intent $\Psi_j$ is registered by user $U_i$ at time $t_0$ and is composed of $N$ actions internally, $\{\Psi_{j, 0}, \Psi_{j, 1}, ... \Psi_{j, N-1}\}$. 

After some time $t_1$ (where $t_1 > t_0$), the time turner is able to find for the contigous range $\{0 \leq k < K; K < (N-1)\}$ the solvers $S_{j, k}$ that are have executed the actions $\Psi_{j, k}$ for a fee of $F_{j, k}$. As such, the cumulative fee for executing the actions $\{\Psi_{j, 0}, \Psi_{j, 1}, ... \Psi_{j, K}\}$ is a sum of solver fees and dApp option premiums: 
$$
F_{Partial} = \mathlarger{\sum}_{k = 0}^{k = K-1}F_{j, k} + \mathlarger{\sum}_{k = 0}^{k = K-1}  \mathlarger{\sum}_{l = 0}^{l = ops(j, k)} P_{j, k, l}
$$

We can denote this partial execution using $E_{j, k} = (S, F_{Partial})$.

At this point, if a different set of solvers $S'_{j, 0}, S'_{j, 1}, ... S'_{j, K}$ are able to execute the actions $\Psi_{j, K+1}, \Psi_{j, K+2}, ... \Psi_{j, N-1}$ for a fee of $F'_{j, K+1}, F'_{j, K+2}, ... F'_{j, N-1}$. Similar to above, the cumulative fee for executing the actions $\{\Psi_{j, 0}, \Psi_{j, 1}, ... \Psi_{j, K-1}\}$ is a sum of solver fees and dApp option premiums:
$$
F'_{Partial} = \mathlarger{\sum}_{k = 0}^{k = K-1}F'_{j, k} + \mathlarger{\sum}_{k = 0}^{k = K-1} \mathlarger{\sum}_{l = 0}^{l = ops(j, k)}P'_{j, k, l}
$$

We can denote this partial execution using $E'_{j, k} = (S', F'_{Partial})$.

\textbf{Execution forking}: We say execution $E'$ is preferred over $E$ if $F'_{Partial} < F_{Partial} - R$ where $R$ is the sum of all partial refund of premium one can get upon surrendering the execution's dApp options, the \textbf{premium refund} described below.

This rule does not stand just for partial execution, but also for the entire execution of the intent.

\subsection{DApps' "Forking Loss" compensation}
The dApps involved in the actions of the intent are also supposed to be compensated for the loss of fee due to execution forking. To understand the necessity of this, consider an action reserving the "Option" from a ETH-USDC Swap DApp for "0.5 ETH to 1500 USDC swap" at a later max time of redemption $t_{redeem}$. Holding this option for an intent is like holding an option in traditional finance.

Since the premium the dApp charges is charged for holding the option till time $t_{redeem}$, the question arises that what happens if the option is confirmed to be not exercised / surrendered before $t_{redeem}$ by the time turner? Option surrendering before redemption time can lead to dApps being able to take early decisions and provide better pricing for execution of other actions and hence, the dApp may choose to pay back some portion of the premium to the user. This is called \textbf{premium refund} and is denoted by $R_{j, k, l}$. This is entirely optional and upto dApps to implement however. There is a valid case where the dApp may not want to refund the premium at all.

\subsection{Intent griefing by solvers}

Every execution forking incurs a cost to the user in form of a "forking loss". This is the attributable to the option holding cost on the fork not chosen. This can be exploited by malicious solvers to grief the user by executing the intent in a way that is not optimal for the user.

More precisely, assume an intent $\Psi_j$ executed without any forking in its lifecycle. This intent's execution can make full use of $\Psi_\text{MAXFEE}$. However if we compare this to a case where the intent had some partial execution but then forked to this execution, the user would have to pay an additional fee of holding the dApp options for the previous forked execution, and part of the $\Psi_\text{MAXFEE}$ is lost in paying for options never redeemed.

The way in which a solver can profit from this mechanism is if it colluded with the dApps such that the dApp options are guaranteed to not be redeemed due to a pre-meditated post-hoc execution forking of an intent post option buy from the dApp. Then, a malicious solver can execute the intent partially, make the intent pay for the dApp's option and later fork the execution into a better path, making the dApp pocket the option premium minus premium refund. This profit, also called \textbf{griefing profit} can then be shared between the solver and the dApp via some internal arrangement.

The mechanism of why this is possible is similar to the case of MEV. This is because the solvers find themselves in a special position to be able to affect the execution of the intent which very few actor apart from them can.

\subsection{Execution Payments}
Each intent $\Psi_j$ composed of a set of actions $A_j = \{\Psi_{j, 0} ... \Psi_{j, n}\}$ has to pay the intent's action executing solvers $S_j = \{S_{j, 0} ... S_{j, n}\}$ and dApps involved in each of the actions a fee for advancing the intent and getting the action delta receipts(ADRs) $\{\Delta_{j, 0} ... \Delta_{j, n}\}$. 

The maximum fee amount that the user is willing to pay for the intent over all the constituent actions is called \textbf{max intent fee} and is denoted by $\Psi_\text{MAXFEE}$. It is specified by the user $U_i$ at the time of intent registration. For simplicity of the argument, let us assume that $\Psi_\text{MAXFEE}$ is denominated in the native currency of the blockchain, for example, ETH for Ethereum and SOL for Solana. 

However, the payment amount can be denominated in any currency that holds value on-chain including ERC-20 tokens, and a combination of ERC-20 and native tokens. We call the two formats of payment \textbf{native fee regime} and \textbf{composite fee regime} respectively.

There exists a possibility that the intent upon execution needs more fee than the maximum fee specified by the user at the time of the registration. In such a case, the user may be given an option to either pay the extra fee or cancel the intent (the default behavior in case of a non-response). As such the time turner exposes a function \code{addMaxIntentFee(j)} to increase the max fee of $\Psi_j$.

\subsection{Intent Registration}

\subsection{Intent halt due to EOW collapse}

\subsection{On Options default by dApps or solvers}